\newpage
\phantomsection
\label{prf:inverse-function-derivative}
\LRAProofFor{cor:inverse-function-derivative}

\begin{remark*}[Return]
\hyperref[cor:inverse-function-derivative]{Return to Corollary}
\end{remark*}

\begin{corollary*}[Derivative of the Inverse Function]
Under the hypotheses of the one-variable Inverse Function Theorem, the local inverse \(g\) satisfies \(g'=1/(f'\circ g)\).
\end{corollary*}

\begin{proof}[Professional Standard Proof]
\LRAProofBodyStart
This is conclusion of the one-variable Inverse Function Theorem: for every \(y\in V\),
\[
  g'(y)=\frac{1}{f'(g(y))}.
\]
If \(f:I\to J\) is globally \(C^1\), bijective, and satisfies \(f'\neq 0\) on \(I\), the same local formula holds at every \(y\in J\). These local identities agree on overlaps, so the global inverse is \(C^1\) and satisfies \((f^{-1})'(y)=1/f'(f^{-1}(y))\).
\end{proof}

\begin{proof}[Detailed Learning Proof]
\LRAProofBodyStart
The theorem already constructs the inverse and proves the derivative formula. The corollary only changes the packaging: the formula is read as a named rule. In the global case, each point of \(I\) has a neighbourhood where the local theorem applies; the local inverses are restrictions of the same global inverse, so the derivative formula patches across all of \(J\).
\end{proof}

\begin{remark*}[Proof structure]
Read the local formula from the inverse theorem, then patch the local formulas for the global statement.
\end{remark*}

\begin{dependencies}
\begin{itemize}
  \item \hyperref[thm:inverse-function-theorem-one-variable]{Theorem: Inverse Function Theorem, One Variable}
\end{itemize}
\end{dependencies}

\clearpage
